\chapter{Conclusiones y trabajo futuro}
\label{cap:conclusiones}

Las conclusiones y las limitaciones se redactarán al cierre del desarrollo, cuando pueda
valorarse el grado de cumplimiento de cada objetivo específico. Las líneas de trabajo futuro
identificadas hasta ahora se recogen a continuación.

% Secciones ocultas hasta el cierre del desarrollo:
% \section{Conclusiones}
% TODO: grado de cumplimiento de cada objetivo específico planteado en la introducción;
% contribuciones efectivamente aportadas; lecciones aprendidas del desarrollo.
% \section{Limitaciones}
% TODO: alcance de la validación, dependencia de los conjuntos de datos empleados,
% aspectos no cubiertos por el prototipo.

\section{Líneas de trabajo futuro}

\begin{itemize}
    \item \textbf{Ampliación de la detección.} Incorporar telemetría de red ya disponible en
    los conjuntos de datos, fuentes de Linux, nuevas familias de ataque y simulaciones propias
    de campañas completas en el laboratorio aislado. Además de ampliar la cobertura, mostraría
    que la arquitectura generaliza: una fuente nueva solo requiere un mapeador al esquema OCSF
    y sus reglas. Esta línea está planificada como ampliación condicionada a que el núcleo del
    trabajo esté cerrado a finales de febrero.
    \item \textbf{Señuelos.} Desplegar credenciales y recursos señuelo en la red real, con una
    tasa de falsos positivos casi nula. Una credencial señuelo empleada tras un volcado de LSASS
    reforzaría el patrón de robo de credenciales. Los señuelos expuestos a Internet muestran
    atacantes oportunistas y no cómo se atacaría una organización concreta, que es lo que
    aporta el análisis de rutas, y los de alta interacción plantean riesgos técnicos y legales.
    \item \textbf{Cadena de suministro de software}: extender el razonamiento sobre rutas de
    ataque a las dependencias del software desplegado.
    \item \textbf{Generación asistida del grafo}: obtener hechos del grafo de ataque a partir
    de informes en lenguaje natural.
    \item \textbf{Validación en producción}: evaluar el sistema sobre telemetría real de una
    organización, con su volumen y su actividad legítima de fondo.
\end{itemize}
